\documentclass{article}%
\usepackage{amsmath}%
\usepackage{amsfonts}%
\usepackage{amssymb}%
\usepackage{graphicx}
\usepackage{sectsty}
\usepackage{listings}
\usepackage{framed}
\usepackage{color}
\usepackage{verbatim}
\usepackage{graphicx} %插入图片的宏包
\usepackage{float} %设置图片浮动位置的宏包
\usepackage{subfigure} %插入多图时用子图显示的宏包
\definecolor{shadecolor}{rgb}{.9, .9, .9}
\sectionfont{\fontsize{20}{20}\selectfont}
%-------------------------------------------
\newtheorem{theorem}{Theorem}
\newtheorem{acknowledgement}[theorem]{Acknowledgement}
\newtheorem{algorithm}[theorem]{Algorithm}
\newtheorem{axiom}[theorem]{Axiom}
\newtheorem{case}[theorem]{Case}
\newtheorem{claim}[theorem]{Claim}
\newtheorem{conclusion}[theorem]{Conclusion}
\newtheorem{condition}[theorem]{Condition}
\newtheorem{conjecture}[theorem]{Conjecture}
\newtheorem{corollary}[theorem]{Corollary}
\newtheorem{criterion}[theorem]{Criterion}
\newtheorem{definition}[theorem]{Definition}
\newtheorem{example}[theorem]{Example}
\newtheorem{exercise}[theorem]{Exercise}
\newtheorem{lemma}[theorem]{Lemma}
\newtheorem{notation}[theorem]{Notation}
\newtheorem{problem}[theorem]{Problem}
\newtheorem{proposition}[theorem]{Proposition}
\newtheorem{remark}[theorem]{Remark}
\newtheorem{solution}[theorem]{Solution}
\newtheorem{summary}[theorem]{Summary}
\newenvironment{proof}[1][Proof]{\textbf{#1.} }{\ \rule{0.5em}{0.5em}}
\setlength{\textwidth}{7.0in}
\setlength{\oddsidemargin}{-0.35in}
\setlength{\topmargin}{-0.5in}
\setlength{\textheight}{30.0in}
\setlength{\parindent}{0.3in}

\newenvironment{code}%
   {\snugshade\verbatim}%
   {\endverbatim\endsnugshade}

\begin{document}

\begin{flushright}
\begin{LARGE}

\textbf{Zhuang Liu \\
SID: 25727277}
\end{LARGE}
\end{flushright}

\begin{center}
\begin{huge}
\textbf{CS 230 Distributed Computer Systems \\
Homework 3} \\
\end{huge}
\end{center}

\section{Six Nested Loop Cases}
\begin{Large}

Here, I used a Java program to simulate the process. I choose \texttt{n} of 3. At first, I defined three matrices: \texttt{a}, \texttt{b} and \texttt{c} as follows:
\begin{code}
    int[][] a = {
            {1, 2, 3},
            {3, 1, 2},
            {2, 3, 1}
    };

    int[][] b = {
            {1, 1, 1},
            {2, 2, 2},
            {3, 3, 3}
    };

    int[][] c = {
            {1, 2, 3},
            {1, 2, 3},
            {1, 2, 3}
    };
\end{code}
Next, we can use a nested loop to print every step in the matrix manipulation. This is how we mimic the \texttt{(i, j, k)} case:
\begin{code}
    for (int i = 0; i < 3; i++) {
        for (int j = 0; j < 3; j++) {
            for (int k = 0; k < 3; k++) {
                c[i][j] = c[i][j] + a[i][k] * b[k][j];
                System.out.println("c" + i + j + " = " + c[i][j]);
            }
        }
    }
\end{code}
\newpage
By changing the order of the loops, we can mimic other cases. For example, this is how we mimic the \texttt{(i, k, j)} case:
\begin{code}
    for (int i = 0; i < 3; i++) {
        for (int k = 0; k < 3; k++) {
            for (int j = 0; j < 3; j++) {
                c[i][j] = c[i][j] + a[i][k] * b[k][j];
                System.out.println("c" + i + j + " = " + c[i][j]);
            }
        }
    }
\end{code}
In the output of the program, we can get the execution order of the process.
\begin{LARGE}
\begin{flushleft}
\textbf{1.1 (i, j, k) Case}
\end{flushleft}
\end{LARGE}

In this case, the execution order is as follows:
\begin{code}
    c00 c00 c00 c01 c01 c01 c02 c02 c02 
    c10 c10 c10 c11 c11 c11 c12 c12 c12
    c20 c20 c20 c21 c21 c21 c22 c22 c22
\end{code}
As we can see from the order, since \texttt{k} is in the most inner loop, every \texttt{c(i, j)} computation will stay on current \texttt{c(i, j)} until \texttt{k} increased from 0 to \texttt{n - 1}. As the \texttt{i} increases from 0 to \texttt{n - 1}, in each \texttt{i}, \texttt{j} increases from 0 to \texttt{n - 1}.
\begin{LARGE}
\begin{flushleft}
\textbf{1.2 (i, k, j) Case}
\end{flushleft}
\end{LARGE}

In this case, the execution order is as follows:
\begin{code}
    c00 c01 c02 c00 c01 c02 c00 c01 c02 
    c10 c11 c12 c10 c11 c12 c10 c11 c12
    c20 c21 c22 c20 c21 c22 c20 c21 c22
\end{code}
We can regard every 3 \texttt{c(i, j)} as one unit, inside every unit, \texttt{j} increases from 0 to \texttt{n - 1}, since \texttt{j} is in the most inner loop. The computation will stay on the unit until \texttt{k} increased from 0 to \texttt{n - 1}, since the \texttt{k} is on the outer loop of \texttt{j}. The \texttt{i} increases every 3 computation unit from 0 to \texttt{n - 1}.
\newpage
\begin{LARGE}
\begin{flushleft}
\textbf{1.3 (j, i, k) Case}
\end{flushleft}
\end{LARGE}

In this case, the execution order is as follows:
\begin{code}
    c00 c00 c00 c10 c10 c10 c20 c20 c20 
    c01 c01 c01 c11 c11 c11 c21 c21 c21
    c02 c02 c02 c12 c12 c12 c22 c22 c22
\end{code}
As we can see from the order, since \texttt{k} is in the most inner loop, every \texttt{c(i, j)} computation will stay on current \texttt{c(i, j)} until \texttt{k} increased from 0 to \texttt{n - 1}. As the \texttt{j} increases from 0 to \texttt{n - 1}, in each \texttt{j}, \texttt{i} increases from 0 to \texttt{n - 1}.

\begin{LARGE}
\begin{flushleft}
\textbf{1.4 (j, k, i) Case}
\end{flushleft}
\end{LARGE}

In this case, the execution order is as follows:
\begin{code}
    c00 c10 c20 c00 c10 c20 c00 c10 c20 
    c01 c11 c21 c01 c11 c21 c01 c11 c21
    c02 c12 c22 c02 c12 c22 c02 c12 c22
\end{code}
We can regard every 3 \texttt{c(i, j)} as one unit, inside every unit, \texttt{i} increases from 0 to \texttt{n - 1}, since \texttt{i} is in the most inner loop. The computation will stay on the unit until \texttt{k} increased from 0 to \texttt{n - 1}, since the \texttt{k} is on the outer loop of \texttt{i}. The \texttt{j} increases every 3 computation unit from 0 to \texttt{n - 1}.
\begin{LARGE}
\begin{flushleft}
\textbf{1.5 (k, i, j) Case}
\end{flushleft}
\end{LARGE}

In this case, the execution order is as follows:
\begin{code}
    c00 c01 c02 c10 c11 c12 c20 c21 c22 
    c00 c01 c02 c10 c11 c12 c20 c21 c22
    c00 c01 c02 c10 c11 c12 c20 c21 c22
\end{code}
We can regard every 9 \texttt{c(i, j)} as a computation unit. When an unit is finished, the \texttt{k} increments, then start a new unit, this process last until \texttt{k} increases from 0 to \texttt{n - 1}. In side every unit, the \texttt{i} increases from 0 to \texttt{n - 1}. In every \texttt{i}, \texttt{j} increases from 0 to \texttt{n - 1}.
\newpage
\begin{LARGE}
\begin{flushleft}
\textbf{1.6 (k, j, i) Case}
\end{flushleft}
\end{LARGE}

In this case, the execution order is as follows:
\begin{code}
    c00 c10 c20 c01 c11 c21 c02 c12 c22 
    c00 c10 c20 c01 c11 c21 c02 c12 c22
    c00 c10 c20 c01 c11 c21 c02 c12 c22
\end{code}
We can regard every 9 \texttt{c(i, j)} as a computation unit. When an unit is finished, the \texttt{k} increments, then start a new unit, this process last until \texttt{k} increases from 0 to \texttt{n - 1}. In side every unit, the \texttt{j} increases from 0 to \texttt{n - 1}. In every \texttt{j}, \texttt{i} increases from 0 to \texttt{n - 1}.

\end{Large}

\section{N-Processor Ring Mapping}
\begin{Large}

\begin{LARGE}
\begin{flushleft}
\textbf{2.1 (i, j, k) Case}
\end{flushleft}
\end{LARGE}
The SPMD program of this case becomes:
\begin{code}
    Do All P(i = 0, n - 1):
        For j = i, do n times:
            For k = 0, n - 1, do:
                c(i, j) = c(i, j) + a(i, k) + b(k, j)
            endfor
        Every Processor i sends its current 
        column of B to processor (i + 1) mod (n)
        endfor
    enddoall
\end{code}
Now let's see the execution order of this program again:
\begin{code}
    c00 c00 c00 c01 c01 c01 c02 c02 c02 
    c10 c10 c10 c11 c11 c11 c12 c12 c12
    c20 c20 c20 c21 c21 c21 c22 c22 c22
\end{code}
We can see that the first row's \texttt{i} is all 0, while the second row's \texttt{i} and third row's \texttt{i} is 1 and 2. It means we can unfold the program on \texttt{i}. Each row is a processor. In each processor, we only need to store \texttt{a(i, 0 ... n - 1)}, \texttt{b(i, 0 ... n - 1)} and \texttt{c(i, 0 ... n - 1)}, since it's paralleled by i, all the \texttt{i} in the processor is the same.
\newpage
\begin{code}
Time = O(n2) Multiplication
Memory = O(n)
Communication = O(n2) Network Cycles
\end{code}
\begin{LARGE}
\begin{flushleft}
\textbf{2.2 (i, k, j) Case}
\end{flushleft}
\end{LARGE}
The SPMD program of this case becomes:
\begin{code}
    Do All P(i = 0, n - 1):
        For k = i, do n times:
            For j = 0, n - 1, do:
                c(i, j) = c(i, j) + a(i, k) + b(k, j)
            endfor
        Every Processor i sends its current 
        column of B to processor (i + 1) mod (n)
        endfor
    enddoall
\end{code}
Now let's see the execution order of this program again:
\begin{code}
    c00 c01 c02 c00 c01 c02 c00 c01 c02 
    c10 c11 c12 c10 c11 c12 c10 c11 c12
    c20 c21 c22 c20 c21 c22 c20 c21 c22
\end{code}
It is similar to the \texttt{(i, j, k)} case, since the very outer loop is also i, so it also can be paralleled by \texttt{i}. Each row is a processor. In each processor, the \texttt{i} remains unchanged in the process. The storage scheme remains the same. However, because the order of \texttt{j} loop and \texttt{k} loop is swapped, the execution order inside the processor will change too.
\begin{code}
Time = O(n2) Multiplication
Memory = O(n)
Communication = O(n2) Network Cycles
\end{code}

\begin{LARGE}

\begin{flushleft}
\textbf{2.3 (j, i, k) Case}
\end{flushleft}
\end{LARGE}
The SPMD program of this case becomes:
\newpage

\begin{code}
    Do All P(j = 0, n - 1):
        For i = j, do n times:
            For k = 0, n - 1, do:
                c(i, j) = c(i, j) + a(i, k) + b(k, j)
            endfor
        Every Processor j sends its current 
        column of A to processor (j + 1) mod (n)
        endfor
    enddoall
\end{code}
Now let's see the execution order of this program again:
\begin{code}
    c00 c00 c00 c10 c10 c10 c20 c20 c20 
    c01 c01 c01 c11 c11 c11 c21 c21 c21
    c02 c02 c02 c12 c12 c12 c22 c22 c22
\end{code}
We can see that, in each row, the \texttt{j} remains the same. For example, in the first row, the \texttt{j}s are all zero. So we can parallel the program by \texttt{j}. We can regard each row as a processor. In each processor, we only store data of \texttt{a(0 ... n - 1, j)}, \texttt{b(0 ... n - 1, j)} and \texttt{c(0 ... n - 1, j)}, since in every processor, \texttt{j} remains the same inside.

\begin{code}
Time = O(n2) Multiplication
Memory = O(n)
Communication = O(n2) Network Cycles
\end{code}

\begin{LARGE}
\begin{flushleft}
\textbf{2.4 (j, k, i) Case}
\end{flushleft}
\end{LARGE}
The SPMD program of this case becomes:

\begin{code}
    Do All P(j = 0, n - 1):
        For k = j, do n times:
            For i = 0, n - 1, do:
                c(i, j) = c(i, j) + a(i, k) + b(k, j)
            endfor
        Every Processor j sends its current 
        column of A to processor (j + 1) mod (n)
        endfor
    enddoall
\end{code}
\newpage
Now let's see the execution order of this program again:
\begin{code}
    c00 c10 c20 c00 c10 c20 c00 c10 c20 
    c01 c11 c21 c01 c11 c21 c01 c11 c21
    c02 c12 c22 c02 c12 c22 c02 c12 c22
\end{code}
This case is similar to the case \texttt{(j, i, k)}. In this case, we can also parallel the process by \texttt{j}. Each row is a processor. Inside each processor, \texttt{j} remains the same. The storage scheme is same with the case \texttt{(j, i, k)}, though the execution order inside the processor differs from the case \texttt{(j, i, k)}.
\begin{code}
Time = O(n2) Multiplication
Memory = O(n)
Communication = O(n2) Network Cycles
\end{code}
\begin{LARGE}
\begin{flushleft}
\textbf{2.5 (k, i, j) Case}
\end{flushleft}
\end{LARGE}
The SPMD program of this case becomes:

\begin{code}
    Do All P(k = 0, n - 1):
        For i = k, do n times:
            For j = 0, n - 1, do:
                c(i, j) = c(i, j) + a(i, k) + b(k, j)
            endfor
        Every Processor k sends its current 
        column of A to processor (i + 1) mod (n)
        endfor
    enddoall
\end{code}
Now let's see the execution order of this program again:
\begin{code}
    c00 c01 c02 c10 c11 c12 c20 c21 c22 
    c00 c01 c02 c10 c11 c12 c20 c21 c22
    c00 c01 c02 c10 c11 c12 c20 c21 c22
\end{code}
In this case, we unfold the program on \texttt{k}. As we can see from the order above, three rows are identical. It is because inside every processor, there is always \texttt{i - j} nested loop. In each processor, we need to store all matrix \texttt{a} and matrix \texttt{b}. However, since the \texttt{k} inside one single processor is unchanged, we don't need to store the whole matrix \texttt{c} inside the processor.
\newpage
\begin{code}
Time = O(n2) Multiplication
Memory = O(n)
Communication = 0 Network Cycles
\end{code}
\begin{LARGE}
\begin{flushleft}
\textbf{2.6 (k, j, i) Case}
\end{flushleft}
\end{LARGE}
The SPMD program of this case becomes:

\begin{code}
    Do All P(k = 0, n - 1):
        For j = k, do n times:
            For i = 0, n - 1, do:
                c(i, j) = c(i, j) + a(i, k) + b(k, j)
            endfor
        Every Processor k sends its current 
        column of B to processor (j + 1) mod (n)
        endfor
    enddoall
\end{code}
Now let's see the execution order of this program again:
\begin{code}
    c00 c10 c20 c01 c11 c21 c02 c12 c22 
    c00 c10 c20 c01 c11 c21 c02 c12 c22
    c00 c10 c20 c01 c11 c21 c02 c12 c22
\end{code}
In this case, we also unfold the program on \texttt{k}. The storage scheme is the same as case \texttt{(k, i, j)}. However, the execution order inside the processor is different from the case \texttt{(k, i, j)}.
\begin{code}
Time = O(n2) Multiplication
Memory = O(n)
Communication = 0 Network Cycles
\end{code}
\end{Large}


\end{document}